% GENERATED FILE. Edit canonical lessons, then run npm run generate:books.
% canonical-source-hash: fnv1a64:78a3d63193fcedfc
% canonical-lessons: SA-C473-angam, SA-C473-jangha, SA-C473-parsnih, SA-C473-mustih, SA-C473-akuncayati

\chapter{Arms and Legs}
\label{ch:sa-arms-and-legs}
\hlchaptermodality{473}

\begin{chapteropening}
\textbf{By the end of this chapter:} \emph{I can name parts of the arms and legs.}

Complete the last lesson of chapter 473: I can name parts of the arms and legs.
\end{chapteropening}

\section[\texorpdfstring{\sk{अंगम्} (aṅgam)}{अंगम् (aṅgam)}]{\sk{अंगम्} (aṅgam) — a limb}

\label{lesson:SA-C473-angam}



\begin{quote}
Before the new one: say the Sanskrit for spits, then the Sanskrit for a taste.
\end{quote}

\subsection*{You'll want to know: \sk{अंगम्}}
\textbf{\sk{अंगम्}} — \emph{aṅgam} — \textquotedblleft{}a limb\textquotedblright{}.

\begin{quote}
\sk{व्यायामेन} \sk{अंगानि} \sk{दृढानि} \sk{भवन्ति।}
\end{quote}

\begin{quote}
\emph{vyāyāmena aṅgāni dṛḍhāni bhavanti.}
\end{quote}

\begin{quote}
Exercise makes the limbs strong.
\end{quote}

A neuter noun.

\begin{cousinweb}
\textbf{\sk{अंगम्}} (\emph{aṅgam}) is traced to a root √am in the Uṇādi list that MW follows; MW gives \textquotedblleft{}a limb of the body, a member\textquotedblright{} (MW).
\end{cousinweb}

\begin{grammarlens}[title={What you've built}]
One more word for describing what you lived and felt.
\end{grammarlens}

\subsection*{Your turn}
\begin{itemize}
\raggedright
  \item \emph{Say it:} \emph{aṅgam}
  \item \emph{Say it:} \emph{aṅgam}, once more
  \item \emph{Say it:} \emph{vyāyāmena aṅgāni dṛḍhāni bhavanti.} — Exercise makes the limbs strong.
  \item \emph{Recall:} say the Sanskrit for spits, then the Sanskrit for a taste, then say \emph{aṅgam} again
\end{itemize}

\begin{tcolorbox}[breakable,colback=teal!4,colframe=teal!35!black,title={Before you move on}]
What does \sk{अंगम्} mean? (\textquotedblleft{}A limb\textquotedblright{}.) Say it once more.
\end{tcolorbox}
\section[\texorpdfstring{\sk{जंघा} (jaṅghā)}{जंघा (jaṅghā)}]{\sk{जंघा} (jaṅghā) — the shin}

\label{lesson:SA-C473-jangha}



\begin{quote}
Before the new one: say the Sanskrit for a taste, then the Sanskrit for a limb.
\end{quote}

\subsection*{You'll want to know: \sk{जंघा}}
\textbf{\sk{जंघा}} — \emph{jaṅghā} — \textquotedblleft{}the shin\textquotedblright{}.

\begin{quote}
\sk{मम} \sk{जंघायां} \sk{पीडा} \sk{अस्ति।}
\end{quote}

\begin{quote}
\emph{mama jaṅghāyāṃ pīḍā asti.}
\end{quote}

\begin{quote}
I have a pain in my shin.
\end{quote}

A feminine noun.

\begin{cousinweb}
\textbf{\sk{जंघा}} (\emph{jaṅghā}) is traced to a root √jaṃh by MW, which gives \textquotedblleft{}the shank, from the ankle to the knee\textquotedblright{} (MW).
\end{cousinweb}

\begin{grammarlens}[title={What you've built}]
One more word for describing what you lived and felt.
\end{grammarlens}

\subsection*{Your turn}
\begin{itemize}
\raggedright
  \item \emph{Say it:} \emph{jaṅghā}
  \item \emph{Say it:} \emph{jaṅghā}, once more
  \item \emph{Say it:} \emph{mama jaṅghāyāṃ pīḍā asti.} — I have a pain in my shin.
  \item \emph{Recall:} say the Sanskrit for a taste, then the Sanskrit for a limb, then say \emph{jaṅghā} again
\end{itemize}

\begin{tcolorbox}[breakable,colback=teal!4,colframe=teal!35!black,title={Before you move on}]
What does \sk{जंघा} mean? (\textquotedblleft{}The shin\textquotedblright{}.) Say it once more.
\end{tcolorbox}
\section[\texorpdfstring{\sk{पार्ष्णिः} (pārṣṇiḥ)}{पार्ष्णिः (pārṣṇiḥ)}]{\sk{पार्ष्णिः} (pārṣṇiḥ) — the heel}

\label{lesson:SA-C473-parsnih}



\begin{quote}
Before the new one: say the Sanskrit for a limb, then the Sanskrit for the shin.
\end{quote}

\subsection*{You'll want to know: \sk{पार्ष्णिः}}
\textbf{\sk{पार्ष्णिः}} — \emph{pārṣṇiḥ} — \textquotedblleft{}the heel\textquotedblright{}.

\begin{quote}
\sk{नवं} \sk{पादत्राणं} \sk{मम} \sk{पार्ष्णिं} \sk{पीडयति।}
\end{quote}

\begin{quote}
\emph{navaṃ pādatrāṇaṃ mama pārṣṇiṃ pīḍayati.}
\end{quote}

\begin{quote}
The new shoe hurts my heel.
\end{quote}

A feminine noun.

\begin{cousinweb}
\textbf{\sk{पार्ष्णिः}} (\emph{pārṣṇiḥ}) MW asks, with a question mark, whether it comes from a root √pṛṣ, so its origin is uncertain; it gives \textquotedblleft{}the heel\textquotedblright{} (MW).
\end{cousinweb}

\begin{grammarlens}[title={What you've built}]
One more word for describing what you lived and felt.
\end{grammarlens}

\subsection*{Your turn}
\begin{itemize}
\raggedright
  \item \emph{Say it:} \emph{pārṣṇiḥ}
  \item \emph{Say it:} \emph{pārṣṇiḥ}, once more
  \item \emph{Say it:} \emph{navaṃ pādatrāṇaṃ mama pārṣṇiṃ pīḍayati.} — The new shoe hurts my heel.
  \item \emph{Recall:} say the Sanskrit for a limb, then the Sanskrit for the shin, then say \emph{pārṣṇiḥ} again
\end{itemize}

\begin{tcolorbox}[breakable,colback=teal!4,colframe=teal!35!black,title={Before you move on}]
What does \sk{पार्ष्णिः} mean? (\textquotedblleft{}The heel\textquotedblright{}.) Say it once more.
\end{tcolorbox}
\section[\texorpdfstring{\sk{मुष्टिः} (muṣṭiḥ)}{मुष्टिः (muṣṭiḥ)}]{\sk{मुष्टिः} (muṣṭiḥ) — a fist}

\label{lesson:SA-C473-mustih}



\begin{quote}
Before the new one: say the Sanskrit for the shin, then the Sanskrit for the heel.
\end{quote}

\subsection*{You'll want to know: \sk{मुष्टिः}}
\textbf{\sk{मुष्टिः}} — \emph{muṣṭiḥ} — \textquotedblleft{}a fist\textquotedblright{}.

\begin{quote}
\sk{कुपितः} \sk{बालः} \sk{मुष्टिं} \sk{बध्नाति।}
\end{quote}

\begin{quote}
\emph{kupitaḥ bālaḥ muṣṭiṃ badhnāti.}
\end{quote}

\begin{quote}
The angry boy clenches his fist.
\end{quote}

A masculine noun.

\begin{cousinweb}
\textbf{\sk{मुष्टिः}} (\emph{muṣṭiḥ}) MW suggests, with a \textquotedblleft{}perhaps\textquotedblright{}, that it first meant the hand closed on something stolen (√muṣ, \textquotedblleft{}to steal\textquotedblright{}); the history is not settled (MW).
\end{cousinweb}

\begin{grammarlens}[title={What you've built}]
One more word for describing what you lived and felt.
\end{grammarlens}

\subsection*{Your turn}
\begin{itemize}
\raggedright
  \item \emph{Say it:} \emph{muṣṭiḥ}
  \item \emph{Say it:} \emph{muṣṭiḥ}, once more
  \item \emph{Say it:} \emph{kupitaḥ bālaḥ muṣṭiṃ badhnāti.} — The angry boy clenches his fist.
  \item \emph{Recall:} say the Sanskrit for the shin, then the Sanskrit for the heel, then say \emph{muṣṭiḥ} again
\end{itemize}

\begin{tcolorbox}[breakable,colback=teal!4,colframe=teal!35!black,title={Before you move on}]
What does \sk{मुष्टिः} mean? (\textquotedblleft{}A fist\textquotedblright{}.) Say it once more.
\end{tcolorbox}
\section[\texorpdfstring{\sk{आकुञ्चयति} (ākuñcayati)}{आकुञ्चयति (ākuñcayati)}]{\sk{आकुञ्चयति} (ākuñcayati) — bends a limb}

\label{lesson:SA-C473-akuncayati}



\begin{quote}
Before the new one: say the Sanskrit for the heel, then the Sanskrit for a fist.
\end{quote}

\subsection*{You'll want to know: \sk{आकुञ्चयति}}
\textbf{\sk{आकुञ्चयति}} — \emph{ākuñcayati} — \textquotedblleft{}bends a limb\textquotedblright{}.

\begin{quote}
\sk{वृद्धः} \sk{शनैः} \sk{जानु} \sk{आकुञ्चयति।}
\end{quote}

\begin{quote}
\emph{vṛddhaḥ śanaiḥ jānu ākuñcayati.}
\end{quote}

\begin{quote}
The old man slowly bends his knee.
\end{quote}

Past (he …): \sk{सः} \sk{जानु} \sk{आकुञ्चितवान्}, \emph{saḥ jānu ākuñcitavān}.

\begin{cousinweb}
\textbf{\sk{आकुञ्चयति}} (\emph{ākuñcayati}) is built from \emph{ā} and the root √kuñc, \textquotedblleft{}to bend\textquotedblright{}; MW gives the causative \textquotedblleft{}to bend (as a limb)\textquotedblright{} and the ppp \emph{ākuñcita} \textquotedblleft{}bent (as the arm or the knee)\textquotedblright{} (MW).
\end{cousinweb}

\begin{grammarlens}[title={What you've built}]
One more word for describing what you lived and felt.
\end{grammarlens}

\subsection*{Your turn}
\begin{itemize}
\raggedright
  \item \emph{Say it:} \emph{ākuñcayati}
  \item \emph{Say it:} \emph{ākuñcayati}, once more
  \item \emph{Say it:} \emph{vṛddhaḥ śanaiḥ jānu ākuñcayati.} — The old man slowly bends his knee.
  \item \emph{Recall:} say the Sanskrit for the heel, then the Sanskrit for a fist, then say \emph{ākuñcayati} again
\end{itemize}

\begin{tcolorbox}[breakable,colback=teal!4,colframe=teal!35!black,title={Before you move on}]
What does \sk{आकुञ्चयति} mean? (\textquotedblleft{}Bends a limb\textquotedblright{}.) Say it once more.
\end{tcolorbox}
